Before delving deeper into the \gls{DNP} group parameters, it is useful to first understand \glspl{DNP} conceptually because they are an important consideration in nuclear reactors \cite{das1993importance}.
\Glspl{DNP} are nuclides which decay via $\beta^-$ decay, producing a daughter, called an emitter, which emits a neutron \cite{bank2000jef}.
%\glspl{DNP} are an important consideration in nuclear reactors \cite{das1993importance}.
%An important consideration in nuclear reactors, especially in \glspl{MSR}, are \glspl{DNP}, discovered in 1939, with initial spectral measurements in 1946 \cite{roberts1939further, das1993importance, burgy1946energies}.
In this work, the term \gls{DNP} will be solely used rather than referring to emitters, as the delay of the neutron emission is effectively only due to the $\beta^{-}$ decay \cite{bohr1939mechanism, ott1985introductory}.
Distinguishing between the emitter and the \gls{DNP} is not relevant for this work.

Equation \eqref{eq:simple-reactor-period} shows a simplified equation for the reactor period $\tau_e$, calculated using a single \gls{DNP} group and without the effect of prompt neutrons included \cite{NRC2006EquationsConversions}.
The delay in neutron emission significantly affects the reactor period, or the time it takes for the reactor power to increase by a factor of $e$, and transient behavior in general.
\begin{equation}
    \tau_e \approx \frac{\beta_{eff} - \rho}{\rho \lambda_{eff}}
    \label{eq:simple-reactor-period}
\end{equation}
where
\begin{align}
    \beta_{eff} &= \text{the effective delayed neutron fraction [-]} \nonumber \\
    \rho &= \text{the reactivity [-]} \nonumber \\
    \lambda_{eff} &= \text{the one \gls{DNP} group decay constant [$s^{-1}$]}. \nonumber
\end{align}
The effective delayed neutron fraction is discussed in more detail in Section \ref{sec:gandkparams}.
The one \gls{DNP} group approximation is not used in this study, so this approximation is not discussed in greater detail.
However, this simplified equation demonstrates how the reactor period depends greatly on the effective delayed neutron fraction, which itself is based on the \glspl{DNP}.

Equation \eqref{eq:dnpsingleeq} shows the time-dependent rate of delayed neutron emission from a single \gls{DNP} as:
\begin{equation}
    \dot{n}_{d, i}(t) = P_{n, i} \lambda_i N_i(t),
    \label{eq:dnpsingleeq}
\end{equation}
\begin{align}
\intertext{where}
\dot{n}_{d, i} &= \mbox{ the delayed neutron activity, or count rate, from \gls{DNP} $i$ $\left[n \cdot s^{-1}\right]$} \nonumber \\
P_{n, i} &= \mbox{ the average number of delayed neutrons emitted from} \nonumber\\
& \;\;\;\;\;\; \mbox{decay of nuclide $i$, or the emission probability [-]} \nonumber \\
\lambda_i &= \mbox{ the decay constant of \gls{DNP} $i$ $[s^{-1}]$} \nonumber \\
N_i(t) &= \mbox{ the number of atoms of \gls{DNP} $i$ $\left[atoms\right]$.} \nonumber
\end{align}

%\begin{itemize}
%    \setlength\itemsep{-0.25em}
%    \item $n_{d, i}$ is the delayed neutron activity, or count rate, from \gls{DNP} $i$
%    \item $P_{n, i}$ is the average number of delayed neutrons emitted from a single beta decay of nuclide $i$, or the emission probability
%    \item $\lambda_i$ is the decay constant of \gls{DNP} $i$
%    \item $N_i(t)$ is the concentration of \gls{DNP} $i$
%\end{itemize}

%A challenging aspect of treating \glspl{DNP} in \glspl{MSR} is the necessary inclusion of spatial resolution due to the movement of the fuel salt.
%In \glspl{MSR}, the movement of the fuel salt necessitates spatial resolution for depletion, transmutation, and neutron emission.
In a solid-fueled reactor, such as a \gls{LWR}, spatial resolution is not necessary to properly model the \glspl{DNP}, as the fuel remains stationary.
In \glspl{MSR}, the movement of the fuel salt necessitates spatial resolution for tracking \gls{DNP} distributions and concentrations.
Modeling space and time together largely increases computational cost, making problems involving \glspl{DNP} in \glspl{MSR} more challenging than handling \glspl{DNP} in stationary fuels.

Properly modeling the \glspl{DNP} in reactors is important, otherwise simulations may predict incorrect power responses, reactivity feedback, control element worth, and reactor controllability \cite{das1993importance, das1977kinetics, lane1958fluid}.
Section \ref{sec:bg-trans-analysis-methods} shows how the power response depends on the \glspl{DNP}.
The reactor controllability varies based on the available number of \glspl{DNP}.
If there are fewer delayed neutrons available due to movement of the fuel salt, such as in an \gls{MSR}, then the reactor period becomes shorter, making the reactor harder to control.